\documentclass[10pt,a4paper]{amsart}
\usepackage[margin=19mm]{geometry}
\usepackage{amsmath,amssymb,mathtools}
\usepackage[scr=rsfs,cal=euler]{mathalfa}
\usepackage{xfrac}
\usepackage[unicode,hidelinks,bookmarks=false]{hyperref}
\newcommand{\ie}{i.e.\ }
\newcommand{\cf}{cf.\ }
\newcommand{\resp}{respectively\ }
\newcommand{\Subsection}{\subsection}
\providecommand{\nobreakdash}{\nobreak\hbox{-}}
\setlength{\emergencystretch}{2em}
\multlinegap0pt
\usepackage{enumerate}
\usepackage{amssymb}
\usepackage[shortlabels]{enumitem}
\usepackage{tikz}
\newtheorem{proposition}{Proposition}
\newtheorem{remark}{Remark}
\newtheorem{corollary}{Corollary}
\newtheorem{theorem}{Theorem}
\newtheorem{lemma}{Lemma}
\newcommand{\E}{\mathbb{E}}
\title{Concentration inequalities for maximal displacement of random walks on groups of polynomial growth}
\author{J\'er\'emie Brieussel, Romain Tessera, Tianyi Zheng}
\date{Source: arXiv:2601.19417v1 (CC BY 4.0)}
\begin{document}
\maketitle

\noindent\emph{Source and licence.} Concentration inequalities for maximal displacement of random walks on groups of polynomial growth. Version arXiv:2601.19417v1 (CC BY 4.0), distributed by arXiv under the Creative Commons Attribution 4.0 International licence, http://creativecommons.org/licenses/by/4.0/, as recorded in the arXiv licence metadata for this exact version and shown on the abstract page https://arxiv.org/abs/2601.19417v1. Full licence notice: \texttt{licenses/arxiv-2601.19417-CC-BY-4.0.txt}.\par\smallskip

\noindent\textit{Editorial scope.} Counted unit: the article's proposition prop:SubexpSubGauss (source0.tex lines 532--543). The article's complete original proof of it is delivered with it (source0.tex lines 1020--1100, one \texttt{proof} environment of 81 lines, which the article places away from the statement). No mathematics is written, repaired or re-derived: the statement and the proof are the article's own lines, in the article's own notation, and no retained line is substituted or re-flowed.  Retained uncounted as the source-local context the proof needs: lines 550--552; lines 565--569; lines 576--584; lines 603--606; lines 917--926; lines 958--960; lines 964--966; lines 968--974; lines 982--984; lines 1008--1014.  Nothing was omitted from any retained span, so the package carries no omission record.  No retained line contains a citation, so the wrapper carries no bibliography and no bibliography entry.  No retained line contains a figure environment, an image-inclusion command or drawing code, so no image is delivered and the package declares no asset dependency.  Editorial formatting only, in addition: an AMS \texttt{amsart} preamble that restates the article's own declarations and macros used by the retained lines, namely the environments \texttt{proposition}, \texttt{remark}, \texttt{corollary}, \texttt{theorem}, \texttt{lemma} on separate counters, together with the standard packages those lines need.  The retained environments print in this document as Proposition 1, Remark 1, Corollary 1, Theorem 1, Lemma 1 in order of appearance, whereas the article prints the same items with the numbering of its own full document.  The complete native source of the article as distributed in the arXiv e-print for this exact version is bundled unchanged as \texttt{sources/193454/source0.tex}.
\begin{remark}\label{rk:lambda}
If a family $f$ of random variables is $(E,K,b)$-subexponential, then for any real parameter $\lambda$, the familly $\lambda f$ is $(|\lambda| E, \lambda^2 K, \frac{b}{|\lambda|})$-subexponential.
\end{remark}

\begin{corollary}\label{cor:convexSubExp} 
Let $f=(f_{i})_{i\in I}$ and $h=(h_{i})_{i\in I}$ be two families
of random variable and let $\eta\in[0,1]$. If $f$ and $h$ are $(E,K,b)$-subexponential,
then so is $\eta f+(1-\eta)h$.
\end{corollary}

\begin{proposition}\label{prop:products} Let $f=(f_{i})_{i\in I}$
and $h=(h_{i})_{i\in I}$ be two families of real random variables.
If $f$ is $\alpha$-concentrated and $g$ is $\beta$-concentrated
for some $\alpha,\beta>0$, then $fh$ is $\gamma$-concentrated for
$\gamma=\frac{\alpha\beta}{\alpha+\beta}$. 

In particular, if both
$f$ and $h$ are subgaussian, then their product is subexponential.
\end{proposition} 

\begin{proposition}\label{prop:MaxConcentrated} Let $0<\alpha<\infty$,
and let $(X_{n})$ be a martingale. Denote $X_{n}^{*}=\max_{1\leq k\leq n}X_{n}$.
Then if $(X_{n})$ is $\alpha$-concentrated, then so is $(X_{n}^{*})$
\end{proposition} 

\begin{theorem}\label{thm:norm} Let $\mu$ be a compactly supported measure on $G=N \rtimes Q$. Denote $w_n$ the random walk at time $n$. Then the family
\[
\left( \frac{1}{\sqrt{n}}\max_{k\le n}\left(\left|w_k\exp\left(-kv_{\mu}\right)\right|_{\mathfrak{F}_\mu}\right) \right)_{n\in \mathbb{N}}
\]
has subgaussian concentration, i.e. there exists $c_1,c_2>0$ such that
\[
\sup_{n\in \mathbb{N}} \mathbb{P} \left( \frac{1}{\sqrt{n}}\max_{k\le n}\left(\left|w_k\exp\left(-kv_{\mu}\right)\right|_{\mathfrak{F}_\mu}\right) \ge t\right) \le c_2e^{-c_1t^2}.
\]
The constants $c_1,c_2$ depend only on the Lie group via its dimension $d=\mathrm{dim}(N)$ and degree of nilpotency $s$, and on the measure $\mu$ via  the radius $R_\mu$, and the spectral constant $\kappa_\mu$ .
\end{theorem}

\begin{align}\label{eq:increment}
y_n=y_{n-1}\ast (n-1)v_\mu \ast q_{n-1}\zeta_n\ast \left( -(n-1)v_\mu\right),
\end{align}

\begin{align}\label{eq:modn4}
v\ast\zeta\ast(-v) = \zeta+[v,\zeta] \mod \mathfrak{n}^{(4)}.
\end{align}

\begin{remark}\label{rk:power2}
It is sufficient to prove Theorem~\ref{thm:norm} in the case $n$ is a power of $2$. Indeed, if $\frac{1}{\sqrt{2n}}\max_{k\le 2n}X_k$ is subgaussian with constants $c_1,c_2$, then for any $n\le p \le 2n$ the variable $\frac{1}{\sqrt{p}}\max_{k\le p}X_k$ is subgaussian with
\[
\mathbb{P}\left(\frac{1}{\sqrt{p}}\max_{k\le p}X_k \ge s\right)\le 
\mathbb{P}\left(\frac{1}{\sqrt{2n}}\max_{k\le 2n}X_k \ge \frac{s}{\sqrt{2}}\right) \le c_2e^{-\frac{c_1}{2}s^2}.
\]
\end{remark}

\begin{align}\label{eq:y2n}
(y_{2n},q_{2n})=(y_n \ast nv_\mu\ast q_ny'_n,\ast(-nv\mu), q_nq'_n).
\end{align}

 \begin{lemma}\label{lem:theta2}
 The family
 \[
\left(  \frac{1}{\sqrt{n}} \max_{k\le n} \|\theta_2(y_k)\|^{\frac{1}{2}} \right)_{n \in \mathbb{N}}
 \]
 has subgaussian concentration, with constants depending only on $d,R_\mu$, and $\kappa_\mu$.
 \end{lemma}

\begin{proposition}\label{prop:SubexpSubGauss} A family $f=(f_{i})_{i\in I}$
of random variables is 
\begin{itemize}
\item subgaussian if and only if $\sup_{i\in I}|\E(f_{i})|<\infty$ and there exists a constant $K$ such that for all $t \in \mathbb{R}$,
\begin{align}\label{eq:concentLaplace}
\sup_{i\in I}\E\left(\exp(t(f-\E(f))\right)\leq\exp(Kt^{2}).
\end{align}
\item subexponential  if and
only $\sup_{i\in I}|\E(f_{i})|<\infty$ and there exist constants
$(K,b)$  such that (\ref{eq:concentLaplace}) holds  for all $|t|<b$.
\end{itemize}
\end{proposition} 

\begin{proof}[Proof of Lemma~\ref{lem:theta2}]
By induction for $c=1$ we assume that 
\begin{align}\label{eq:theta1}
\left( \frac{1}{\sqrt{n}}\max_{k\le n}\|\theta_1(y_k)\| \right)_n \quad  \textrm{ is subgaussian.}
\end{align}

Equivalently to the lemma, we prove that the family $\left(\frac{1}{n}\max_{k\le n}\|\theta_2(y_n)\|\right)_n$ is subexponential. According to Proposition~\ref{prop:SubexpSubGauss}, we have to show that there exist $E,K,b>0$ such that for all unit vector $u \in \mathfrak{m}^{(2)}$, 
\begin{align}\label{eq:scalu}
\left(\frac{1}{n} \max_{k \le n} \langle \theta_2(y_k) |u \rangle \right)_n \quad \textrm{ is }E,K,b\textrm{-subexponential.}
\end{align}

First (\ref{eq:increment}), (\ref{eq:modn4}) and BCH formula give
\[
\theta_2(y_n)=\theta_2(y_{n-1})+\theta_2(q_{n-1}\zeta_n)+\frac{1}{2}\left[\theta_1(y_{n-1}),\theta_1(q_{n-1}\zeta_n)\right]
\]
by induction.
As $\zeta_n$ is centred, we have $\mathbb{E}\left[\theta_1(y_{n-1}),\theta_1(q_{n-1}\zeta_n)\right]=0$. It follows that 
\[
\mathbb{E}\theta_2(y_n)=\mathbb{E}\theta_2(y_{n-1})+\mathbb{E}\theta_2(q_{n-1}\zeta_n)=\sum_{k=1}^n \mathbb{E}\theta_2(q_{k-1}\zeta_k).
\]
We get $\frac{1}{n}\mathbb{E}\theta_2(y_n) \le E=\sup_{q,\zeta}\|\theta_2(q\zeta)\|=\sup_\zeta\|\theta_2(\zeta)\|$. Here the supremum is taken over $q\in Q$ and $\zeta \in \mathrm{supp}(\pi^\mathfrak{n}\mu)-v_\mu$. It does not depend on $Q$ as its action preserves the norm and the filtration.  
The set $\mathrm{supp}(\pi^\mathfrak{n}\mu)-v_\mu$ is compact, bounded in terms of $R_\mu$ and~$\kappa_\mu$.

Given a vector valued random variable $Z$, we denote $Z^0=Z-\E(Z)$ the recentred random variable. There remains to show the existence of $K,b>0$ such that $\left(\frac{1}{n}\theta_2(y_n)^0\right)_n$ is $(K,b)$-subexponential. According to Remark~\ref{rk:power2}, it is sufficient to do it for $n=2^k$ which we do by induction on $k$. 

Now by (\ref{eq:y2n}), (\ref{eq:modn4}) and the BCH formula, we have
\begin{align*}
\theta_2(y_{2n})=\theta_2(y_n)+\theta_2(q_ny'_n)+\frac{1}{2}\left[ \theta_1(y_n),\theta_1(q_n y'_n)\right].
\end{align*}
As $\theta_1(y_n)$ is centred, we get
\begin{align}\label{eq:th2y}
\theta_2(y_{2n})^0=\theta_2(y_n)^0+\theta_2(q_ny'_n)^0+\frac{1}{2}\left[ \theta_1(y_n),\theta_1(q_n y'_n)\right].
\end{align}

As we chose a norm with bilinearity constant $C_{\|\cdot\|} \le 1$, we have
\[
\left\|\frac{1}{2}\left[ \theta_1(y_n),\theta_1(q_n y'_n)\right]\right\| \le \frac{1}{2} \left\| \theta_1(y_n)\right\| . \left\| \theta_1(q_ny'_n)\right\|.
\]

By assumption (\ref{eq:theta1}), the sequence $\frac{1}{\sqrt{n}}\theta_1(y_n)$, and hence the sequence $\frac{1}{\sqrt{n}}\theta_1(q_ny'_n)$ as well, is subgaussian, so their product is subexponential. By Proposition~\ref{prop:products}, this gives $K_1,b_1>0$, depending on $d,R_\mu$, such that
\begin{align}\label{eq:K1b1}
\left( \frac{1}{2n}\left[ \theta_1(y_n),\theta_1(q_n y'_n)\right] \right)_n \textrm{ is }(K_1,b_1)-\textrm{subexponential.}
\end{align}

Consider $K\ge K_1/4(1-\frac{1}{\sqrt{2}})^2$ and $b < b_1 2(1-\frac{1}{\sqrt{2}})$. We prove by induction that when $n=2^k$, the random variable $\frac{1}{n} \theta_2(y_n)^0$ is $(K,b)$-subexponential.
(Initiation is done by choosing also $K$ large and $b$ small enough that it holds for $\theta_2(\zeta)$, which has finite support, depending on $R_\mu$.)

Assume now that this is true for $n$ and consider (\ref{eq:th2y}). To ease notations denote $\alpha_n:=\frac{1}{n}\theta_2(y_n)^0$ and $\beta_n:= \frac{1}{n}\theta_2(q_ny'_n)^0$. We have for all $|t|\le b$ and all unit $u \in \mathfrak{m}_{2}$:
\[
\mathbb{E}\exp\left(t \langle \alpha_n ,u\rangle\right) \le \exp(Kt^2)
\]
and also, as the $Q$ adjoint action on $\mathfrak{n}$ preserves the norm, conditioning by the values of the $n$ first increments $g_1,\dots,g_n$, we have
\[
\mathbb{E}\left[\exp\left(t \langle \beta_n ,u\rangle\right)|g_1,\dots,g_n\right]=\mathbb{E}\left[\exp\left(t\left\langle\frac{1}{n}\theta_2(y'_n)^0,q_n u\right\rangle\right)|g_1,\dots,g_n\right] \le \exp(Kt^2).
\]
%Here $\mathcal{B}'_n$ denotes the $\sigma$-algebra generated by the random variables $g_{n+1}, \dots, g_{2n}$.


Our key observation is that being ``essentially'' independent, $\alpha_n$ and $\beta_n$ are such that 
$\frac{1}{\sqrt{2}}(\alpha_n+\beta_n)$ is $(K,b)$-subexponential. 
Let us check this point. For all $|t|\leq b$
\begin{eqnarray*}
\E\left[\exp\left(t\left\langle\frac{\alpha_n+\beta_n}{\sqrt{2}},u\right\rangle\right) \mid g_1,\dots,g_n\right]& = & \exp\left(t\left\langle\frac{\alpha_{n}}{\sqrt{2}},u\right\rangle\right) \E\left[\exp\left(t\left\langle\frac{\beta_{n}}{\sqrt{2}},u\right\rangle\right) \mid g_1,\dots,g_n\right]\\
                                                            & \leq &  \exp\left(t\left\langle\frac{1}{\sqrt{2}}\alpha_{n},u\right\rangle\right)\exp(Kt^2/2)\\
\end{eqnarray*}
Hence we deduce that 
\[\E\exp\left(t\left\langle\frac{\alpha_n+\beta_n}{\sqrt{2}},u\right\rangle\right)\leq \exp(Kt^2/2)\exp(Kt^2/2)=\exp(Kt^2).\]

We are going to exploit the fact that 
\[
\frac{1}{2}(\alpha_n+\beta_n)=\frac{1}{\sqrt{2}}\left(\frac{1}{\sqrt{2}}(\alpha_n+\beta_n)\right).
\]
To lighten notation, let $\eta=\frac{1}{\sqrt{2}}$. Equation (\ref{eq:th2y}) can be rewritten as
 \[
 \frac{1}{2n}\theta_2(y_{2n})^0= \eta\left(\frac{1}{\sqrt{2}}(\alpha_n+\beta_n)\right)+(1-\eta)\left(\frac{1}{4(1-\eta)n}\left[\theta_1({y}_{n}),\theta_1(q_{n}y'_n)\right]\right).
 \]
 
By (\ref{eq:K1b1}) and Remark~\ref{rk:lambda}, we have that $\frac{1}{4(1-\eta)n}\left[\theta_1({y}_{n}),\theta_1(q_{n}y'_n)\right]$ is $(\frac{K_1}{4(1-\eta)^2},2(1-\eta)b_1)$-subexponential. The choice of $K,b$ and Corollary \ref{cor:convexSubExp} imply that $\frac{1}{2n}\theta_2(y_{2n})^0$ is $(K,b)$-subexponential, finishing the induction.

We have proved that $\left(\frac{1}{2^k}\theta_2(y_{2^k})^0\right)_n$ is subexponential. As $\theta_2(y_n)^0$ is a martingale, Proposition~\ref{prop:MaxConcentrated} ensures that $\left(\frac{1}{2^k}\max_{j\le 2^k}\theta_2(y_j)^0\right)_n$ is subexponential. This proves the lemma for $n$ a power of $2$, which is sufficient by Remark~\ref{rk:power2}. Proposition~\ref{prop:MaxConcentrated} and Remark~\ref{rk:power2} modify the constants of concentration only by universal multiplicative factors.
\end{proof}

\end{document}
